\documentclass[a4paper,10pt]{scrartcl}

\usepackage[ngerman]{babel}
\usepackage[utf8]{inputenc}
\usepackage[T1]{fontenc}
\usepackage[babel,german=quotes]{csquotes}
\usepackage{hyperref}
\usepackage{geometry}
\usepackage{graphicx}
\usepackage{enumitem}
\usepackage{multicol}

\setlength\parindent{0pt}
\setitemize{itemsep=0pt}


\geometry{
	a4paper,
	left=20mm,
	top=20mm,
	bottom=15mm,
	footskip=5mm
}

\subject{Gruppe \textit{Nummer}}
\title{Pflichtenheft für das Thema \glqq\textit{Klick-Tool}\grqq}
\subtitle{Betreuer: \textit{Prof. Dr. Andreas Nüchter}}
\author{\textit{Tom Lerzer, Fabia Frank, Luis Ángel Álvarez Gil}}
\date{\textit{SS2026}}

\begin{document}

\maketitle



\section{Zielbestimmung}

\subsection{Muss-Kriterien}

\begin{itemize}
    \item Laden von Bilddaten in jpg und png Format
    \item Internes Erstellen von png-Dateien bei Eingabe von 3DTK-spezifischen 3D-Dateiformaten
    \item Implementieren eines GUI zur erleichterten Ausführung der implementierten Funktionalitäten
    \item Platzieren von maximal zwei der geladenen Bilder nebeneinander
    \item Zoom-Funktion für X-fachen Zoom jedes geladenen Bildes
    \item Anwählen einzelner Koordinaten in einem Bild mit Subpixelgenauigkeit von mindestens 0,5 Pixel
    \item Anwählen korrespondierender Koordinaten in zwei Bildern mit Subpixelgenauigkeit von mindestens 0,5 Pixel
    \item Visualisierung gewählter Koordinaten
    \item Visualisierung der Korrespondenz gewählter Koordinaten
    \item Ausgabe gewählter 2D-Koordinaten in Form $(x_1, y_1), ..., (x_n, y_n)$ für einfache bzw. $(x_{1.1}, y_{1.1}, y_{1.2}, y_{1.2}),$ $ ..., (x_{n.1}, y_{n.1}, x{n.2}, y_{n.2})$in .txt-Dateiformat
    \item Integration der implementierten Funktionalitäten in 3DTK-Toolkit 
    \item Implementieren von Kommandozeilenparametern sowie deren Beschreibungen für das Aufrufen der implementierten Funktionalitäten
    \item Erstellen eines Entwicklerhandbuchs für die implementierten Funktionalitäten
    \item Erstellen eines Benutzerhandbuchs für die implementierten Funktionalitäten
\end{itemize}

\subsubsection{Halbzeit-Ziele}
\begin{itemize}
\item Laden von Bilddaten in jpg und png Format
\item Internes Erstellen von png-Dateien bei Eingabe von 3DTK-spezifischen 3D-Dateiformaten
\item Implementieren eines GUI zur erleichterten Ausführung der implementierten Funktionalitäten
\item Platzieren von maximal zwei der geladenen Bilder nebeneinander
\item Zoom-Funktion für X-fachen Zoom jedes geladenen Bildes
\item Anwählen einzelner Koordinaten in einem Bild mit Subpixelgenauigkeit von mindestens 0,5 Pixel
\item Anwählen korrespondierender Koordinaten in zwei Bildern mit Subpixelgenauigkeit von mindestens 0,5 Pixel
    
\end{itemize}
\newpage

\subsubsection{Zuständigkeiten}
\begin{itemize}
\item Erstellen der grafischen Nutzeroberfläche für Platzierung der Bilder sowie Visualisierung gewählter Koordinaten:
\item Funktionalität zum Laden der Bilddateien und Umwandeln von 3DTK-spezifischen 3D-Formaten: 
\item Implementieren der Zoom-Funktion: 
\item Funktionalität zum Anwählen und Ausgeben von Koordinaten bzw. korrespondierenden Koordinaten: 
\item Integration der implementierten Funktionalitäten in 3DTK: 
\item Implementieren der Kommandozeilenparameter incl. Beschreibung:
\item Erstellen Entwicklerhandbuch: 
\item Erstellen Benutzerhandbuch: 

\end{itemize}

\subsection{Kann-Kriterien}

\begin{itemize}
    \item Automatische Wahl der optimalen Anordnung der Bilder (nebeneinander oder übereinander) aufgrund Bildgröße und Bildschirmauflösung
    \item Ausblenden einzelner Korrespondenzlinien zwischen Punkten
    \item Anzeige Koordinaten der gewählten Punkte
    \item Undo-Funktion für ausgewählte Punkte
    \item Frei positionierbare Lupenansicht
\end{itemize}

\subsection{Abgrenzungskriterien}

\begin{itemize}
    \item Anwahl korrespondierender Punkte wird ausschließlich manuell durchgeführt
\end{itemize}


\section{Einsatz}

\subsection{Anwendungsbereiche}
\begin{itemize}
	\item Bildverarbeitung
	\item Robotik
	\item Photogrammetrie
\end{itemize}


\subsection{Zielgruppe}

Das Projekt soll durch Studierende und Lehrende der Universität Würzburg im Rahmen der Robotik-Forschung, der Bildverarbeitung sowie allgemein im Bereich Informatik zum Einsatz gebracht werden. Für weitere interessierte Personen ist das Projekt auf GitHub zugänglich.


\section{Umgebung}

\subsection{Software}

\begin{itemize}
	\item Programmiersprache: C++
	\item Bildverarbeitung: OpenCV \textbf{Version XXX, Link}
	\item Nutzeroberfläche: Dear ImGUI \textbf{Version XXX, Link}
	\item 3D-Framework: 3DTK
	\item Entwicklungsumgebung: Visual Studio Code
	\item Compiler: g++
	\item Build-System: CMake, mindestens Version 3.4
	\item Versionsverwaltung: git, Version 2.51
	\item Fensterverwaltung: GLFW, Version 3.4
	\item Grafik-API: OpenGL
\end{itemize}

\subsection{Hardware}

Die Hardware muss alle Anforderungen der im vorherigen Punkt genannten Software erfüllen.

\section{Funktionalität}

\begin{itemize}
	\item Nutzer starten das Programm via Terminal. Beim Start des Programms werden ein oder zwei Parameter in den unter \textit{Daten} angegebenen Formaten eingegeben
	\item Auf dem angezeigten Bild bzw. den angezeigten Bildern kann vor der weiteren Bearbeitung bis zu \textbf{X-Mal} hineingezoomt werden
	\item Nutzer wählen beliebig viele einzelne oder korrespondierende Punkte aus
	\item Programm zeigt die angewählten Punkte und deren Korrespondenzen im GUI an
	\item Durch Anwählen des Buttons \textbf{Ausgabe} initiieren Nutzer die Ausgabe der gewählten Punkte in .txt-Format
\end{itemize}


\section{Daten}\begin{itemize}
	\item \textbf{Eingabe:} Dateien im png oder jpg-Format sowie 3DTK-spezifische Formate, die durch Anwendung der \textit{scan\_to\_panorama}-Funktion in anzeigbare png-Dateien umgewandelt werden
	\item \textbf{Ausgabe:} .txt-Datei mit Koordinaten in Form $(x_1, y_1), ..., (x_n, y_n)$ für einfache bzw. $(x_{1.1}, y_{1.1},$ $  y_{1.2}, y_{1.2}),..., (x_{n.1}, y_{n.1}, x{n.2}, y_{n.2})$ für korrespondierende Punkte
	\item \textbf{3DTK-spezifische Formate:}	rexp, uos, uos\_map, uos\_rgb, uos\_frames, uos\_map\_frames, old, rts, 	rts\_map, ifp, riegl\_txt, riegl\_rgb, riegl\_bin, zahn, ply, las
	
\end{itemize}



\section{Leistungen}

Das Laden der Bilddateien, die Umwandlung der 3DTK-spezifischen Datenformate, die Zoom-Funktionalität sowie Auswahl und Ausgabe der gewählten Koordinaten soll weitestgehend performant sein und so in einer angemessenen Zeit geschehen.


\section{Benutzungsoberfläche}

\textit{Dieser Abschnitt kann bei Themen ohne GUI weggelassen werden; ansonsten grobe Skizzen der wichtigsten GUI-Ansichten mit jeweils kurzer Beschreibung (Stichpukte oder 1 bis 2 Sätze pro Skizze)}


\section{Qualitätsziele}

Der Code soll gut kommentiert und klar strukturiert sein, sodass sich neue Entwicklerinnen und Entwickler gut in dem Projekt zurechtfinden und es leicht erweitern und ändern können. Um dies zu erleichtern, wird außerdem ein Developerhandbuch erstellt.\\
Das Ergebnis des Projekts soll eine hohe Benutzerfreundlichkeit haben und für Anwenderinnen und Anwender wird zusätzlich ein Benutzerhandbuch bereitgestellt.

\end{document}